% Generated by results/thesis_tables.py - do not edit by hand.
\begin{table}[H]
\centering
\small
\begin{tabular}{|l|r|r|r|r|r|}
\hline
Διαμόρφωση & TP & Διάμεση (TP) & FP & Διάμεση (FP) & AUROC \\
\hline
Πλήρες σύστημα & 404 & 0,95 & 8 & 0,83 & 0,90 \\
Μόνο δείγματα & 402 & 0,96 & 4 & 0,88 & 0,89 \\
Μόνο πολιτικές & 400 & 0,95 & 6 & 0,65 & 1,00 \\
Χωρίς RAG & 400 & 0,95 & 1 & 0,90 & -- \\
\hline
\end{tabular}
\caption{Αυτο-αναφερόμενη βεβαιότητα (\texttt{confidence}) του μοντέλου στα πακέτα που χαρακτηρίστηκαν \texttt{MALICIOUS}: η μεγαλύτερη τιμή μεταξύ των κακόβουλων Behavior Chains του πακέτου. AUROC: η πιθανότητα ένα τυχαίο σωστό θετικό να έχει μεγαλύτερη βεβαιότητα από ένα τυχαίο false positive (υπολογίζεται για τουλάχιστον δύο false positives).}
\label{tab:confidence}
\end{table}
